% ---
% id: protocol-compiler/ass:spec
% kind: ass
% title: "Specification faithful to the apparatus"
% summary: >-
%   Declared dimensions, alphabets, tables and measurements describe the
%   apparatus: the receiver's in every round whatever the others did and,
%   where no source certificate is demanded, Alice's of two or more letters;
%   every input yields a key. The compiler checks only the last, in part.
% prominence: standard
% status: accepted
% uses: [def:spec]
% assumes: []
% refs: []
% provenance: newly-derived
% short_id: A-SPEC
% verification: never
% category: protocol-structure
% affects: [secrecy]
% user_statement: >-
%   Your declared numbers describe your apparatus: each system has its
%   declared dimension, every alphabet and table (announcements, sifting, the
%   raw key bit) is what your devices and software do, and the key rule
%   yields a key value, never "no key", on every input. Each measuring
%   station, of either party, performs exactly its declared measurement in
%   every round, on that round's input alone; afterpulsing or a dead time
%   violates this. A declared receiver instrument's measurement is the model
%   the engine emits from it; the emission step is recorded, not bounded.
%   Where no source-replacement certificate is demanded (an untrusted source,
%   or a one-dimensional A), Alice's declared measurement of her letter is
%   the one she performs; [[A-ALICE-MEAS]] adds every round, alone. The
%   engine checks that the declaration is well formed, not that it matches
%   your hardware.
% lean:
% ---

The compiler decides, on the declared numbers, every condition of a
\tref{well-formed specification}. What no function of the
\tref{declared data} decides is whether those data are the apparatus: a
mis-declared dimension, or a measurement that is complete as written but is
not the one performed, passes every check. Where a
\tref{source-replacement certificate} is demanded, in
\tref{trusted-source mode} at $d_A\ge2$, Alice's measurement is tied to the
source through that certificate, and clause~(i) says nothing of it. On the
receiver side clause~(i) includes that the measurement keeps no memory from
round to round, on every route, a directly declared family and a chain from
elsewhere included: no \tref{declared datum} describes more than one round. The
key rule's codomain is decided only in part: check V5 refuses a key rule that
withholds a key on part of a passing class, and not one that withholds it on
the whole of such a class.

\begin{assumption}\label{ass:spec}
For a \tref{protocol specification} $\Spec$ (Definition~\ref{def:spec}):
(i) the declared dimensions, alphabets and tables of $\Spec$ describe the
apparatus, and so do its declared measurements, namely: the physical
receiver family $\{\Fph_w\}_{w\in\Vph}$ (Definition~\ref{def:spec}\,(b); in
the general format, each \tref{measured factor}'s family of (b$'$)) is the
measurement performed, in every round, on that round's physical receiver
system alone, whatever happens in the other rounds, their outcomes and public
messages included; and
where no source-replacement certificate is demanded of $\Spec$ (in
\tref{untrusted-source mode}, or at $d_A=1$), $\{M_x\}$
(Definition~\ref{def:spec}\,(a)) is the measurement Alice performs on $A$;
(ii) the key rule $r^{0}_{\hold}$ is valued in $\Zval$.
\end{assumption}
